\origpage{114}{ROBERTSON AND SEYMOUR}
We have been unable to decide if this property is again equivalent to planarity.

\sectionhead{\textsc{References}}

1. H. Bergmann, Ein Planarit\"atskriterium f\"ur endliche Graphen, \textit{Elem. Math.} \textbf{37} (1982), 49--51. 2. P. Buneman, A characterization of rigid circuit graphs, \textit{Discrete Math.} \textbf{9} (1974), 205--212. 3. G. A. Dirac, On rigid circuit graphs, \textit{Abh. Math. Sem. Univ. Hamburg} \textbf{25} (1961), 71--76. 4. P. Erd\"os and L. P\'osa, On independent circuits contained in a graph, \textit{Canad. J. Math.} \textbf{17} (1965), 347--352. 5. F. Gavril, The intersection graphs of subtrees in trees are exactly the chordal graphs, \textit{J. Combin. Theory Ser. B} \textbf{16} (1974), 47--56. 6. A. Gy\'arf\'as and J. Lehel, A Helly-type problem in trees, in ``Combinatorial Theory and Its Applications II,'' (Proc. Colloq. Balatonfured, 1969), pp. 571--584. 7. E. Gy\"ori, On the division of graphs to connected subgraphs, \textit{Combinatorics, Colloq. Math. Soc. J\'anos Bolyai} \textbf{18} (1978), 485--494. 8. L. Lov\'asz, A homology theory for spanning trees of a graph, \textit{Acta Math. Acad. Sci. Hungar.} \textbf{30} (1977), 241--251. 9. Neil Robertson and P. D. Seymour, Graph minors. I. Excluding a forest, \textit{J. Combin. Theory Ser. B} \textbf{35} (1983), 39--61. 10. Neil Robertson and P. D. Seymour, Graph minors. II. Algorithmic aspects of tree-width, \textit{J. Algorithms}, in press. 11. Neil Robertson and P. D. Seymour, Graph minors. IV. Tree-width and well-quasi-ordering, submitted for publication. 12. Neil Robertson and P. D. Seymour, Graph minors. VII. Disjoint paths on a surface, submitted for publication. 13. W. T. Tutte, From matrices to graphs, \textit{Canad. J. Math.} \textbf{16} (1964), 108--127.
